\documentclass[12pt]{article}
\usepackage[margin=1in]{geometry}
\usepackage{amsmath}
\usepackage{booktabs}

\title{\textbf{Problem Set 9:} \\
LASSO and Ridge Regression}
\author{Jhinuk Banerjee}
\date{April 14, 2026}

\begin{document}
\maketitle

\section*{Question 7: Dimensions of Training Data}

The original Housing dataset has 13 predictor variables (excluding \texttt{medv}).
After taking log of the outcome, it converts \texttt{chas} to a
factor and creates all pairwise interactions among the continuous predictors, and generates
6th-degree polynomials of each continuous variable. This prepped the training data with
379 rows and 74 predictor columns (75 total including \texttt{medv}).
This represents 61 additional $X$ variables compared to the original 13 predictors.

\section*{Question 8: LASSO Results}

Using 6-fold cross-validation with $\alpha = 1$ (LASSO), the optimal penalty
parameter selected by cross-validation is $\hat{\lambda} = 0.0099$.
The in-sample RMSE is 0.1558 and the out-of-sample RMSE on the test set
is 0.1711.

\section*{Question 9: Ridge Regression Results}

Using 6-fold cross-validation with $\alpha = 0$ (ridge regression), the optimal
penalty parameter is $\hat{\lambda} = 0.0491$.
The out-of-sample RMSE on the test set is 0.1760.

\section*{Question 10: Discussion}

\subsection*{Can OLS be estimated when columns exceed rows?}

No. Ordinary least squares requires inverting the matrix $X'X$, which becomes
singular when the number of predictors $p$ exceeds the number of observations $N$.
In such a setting, OLS has no unique solution. This is precisely the motivation
for penalized regression methods like LASSO and ridge regression: by adding a
regularization penalty, the optimization problem becomes well-posed even when
$p > N$. In our case, after the recipe the training data has 74 predictors and
379 observations, so OLS would still technically be feasible, but with heavy
overfitting risk.

\subsection*{Bias-Variance Tradeoff}

Both LASSO and ridge regression address the bias-variance tradeoff by shrinking
coefficients toward zero, reducing variance at the cost of introducing some bias.
LASSO ($\alpha = 1$) performs variable selection by shrinking some coefficients
exactly to zero, producing a sparse model. Ridge ($\alpha = 0$) shrinks all
coefficients proportionally but retains all predictors.

In terms of out-of-sample RMSE, LASSO achieves 0.1711 while ridge achieves
0.1760. LASSO slightly outperforms ridge here, suggesting that some of the
engineered features (interactions and polynomial terms) are not genuinely
useful for prediction, and LASSO's hard zeroing of those coefficients reduces
overfitting. Both models perform similarly, indicating neither extreme of the
bias-variance tradeoff dominates strongly in this dataset.

\end{document}